% !TEX root = main.tex
% Shared look of the TikZ diagrams (<chapter>/figures/*.tex): font, palette, text commands
% and styles. A figure only sets its grid (named lengths and coordinates) and its font sizes, and
% draws with these styles; `block` and `instances` read the grid's \boxWidth, \boxHeight and
% \groupMargin.

\usetikzlibrary{arrows.meta, backgrounds, calc, fit, positioning, shadows, shapes.geometric}

% Helvetica, as the standalone diagrams in docs/diagrams/ (helvet, scaled).
\newfontfamily\diagramfont{TeX Gyre Heros}[Scale=0.95, Ligatures=TeX]

% Tailwind palette: one hue per stage of the jack monitor (600 for frames and titles, 50 for the
% group fill, 100 for mode chips), slate for text and lines, the rails in the web interface's red
% and blue. Unplugging is the only colored edge.
\definecolor{detect}{HTML}{0284C7}   \definecolor{detectFill}{HTML}{F0F9FF}   \definecolor{detectTint}{HTML}{E0F2FE}
\definecolor{identify}{HTML}{7C3AED} \definecolor{identifyFill}{HTML}{F5F3FF} \definecolor{identifyTint}{HTML}{EDE9FE}
\definecolor{follow}{HTML}{059669}   \definecolor{followFill}{HTML}{ECFDF5}   \definecolor{followTint}{HTML}{D1FAE5}
\definecolor{output}{HTML}{D97706}   \definecolor{outputFill}{HTML}{FFFBEB}   \definecolor{outputTint}{HTML}{FEF3C7}
\definecolor{other}{HTML}{64748B}    \definecolor{otherFill}{HTML}{F8FAFC}    \definecolor{otherTint}{HTML}{F1F5F9}
\definecolor{unplug}{HTML}{E11D48}
\definecolor{ink}{HTML}{1E293B}
\definecolor{muted}{HTML}{475569}
\definecolor{wire}{HTML}{334155}
\definecolor{railHigh}{HTML}{E11D48}
\definecolor{railLow}{HTML}{0284C7}

% A state's name, above its description.
\NewDocumentCommand\statename{m}{{\diagramNameFont\bfseries\color{ink}#1}}
% A contact pulled up or down through its pull resistor: amssymb's \blacktriangle and
% \blacktriangledown, taken from its font (Latin Modern Math has no such glyphs).
\NewDocumentCommand\pullup{m}{#1\kern0.5pt{\color{railHigh}\usefont{U}{msa}{m}{n}\char"4E}}
\NewDocumentCommand\pulldown{m}{#1\kern0.5pt{\color{railLow}\usefont{U}{msa}{m}{n}\char"48}}


\tikzset{
  % Font sizes, set per figure: group titles, state names, everything else.
  diagram title font/.store in=\diagramTitleFont,
  diagram name font/.store in=\diagramNameFont,
  diagram text font/.store in=\diagramTextFont,
  diagram title font=\footnotesize,
  diagram name font=\scriptsize,
  diagram text font=\scriptsize,
  diagram/.style={
    execute at begin picture={\diagramfont\hyphenpenalty=10000\exhyphenpenalty=10000\relax},
    font=\diagramTextFont,
    every node/.append style={text=ink},
    >={Stealth[length=4pt, width=3.2pt]},
  },
  % Corner radius of every frame, box and backdrop, and of the rail tags.
  corner radius/.store in=\diagramCornerRadius,
  corner radius=2pt,
  tag corner radius/.store in=\diagramTagCornerRadius,
  tag corner radius=1pt,
  % Place a node at a column's x and a row's y (two coordinates).
  cell/.style 2 args={at={(#1 |- #2)}},
  % Frames: a stage's group and a group inside it.
  group/.style={draw=#1, fill=#1Fill, line width=0.8pt, inner sep=0pt,
    rounded corners=\diagramCornerRadius,
    /utils/exec={\def\diagramGroupHue{#1}}},
  subgroup/.style={draw=#1, fill=white!55!#1Fill, line width=0.6pt,
    inner sep=0pt, rounded corners=\diagramCornerRadius},
  % A group that exists more than once, e.g. `group=detect, instances=4` (after `group`): a copy of
  % its frame behind it for each further one, each `stack offset` further up and right, and the
  % count on its top right corner, inset by the figure's \groupMargin. One instance draws the plain
  % group.
  stack offset/.store in=\diagramStackOffset,
  stack offset=1mm,
  instances/.code={%
    \ifnum#1>1
      \pgfmathtruncatemacro\diagramGroupCopies{#1 - 1}%
      \edef\diagramGroupStack{{\diagramGroupHue}{\diagramGroupCopies}{#1}}%
      \pgfkeysalso{group stack/.expand once=\diagramGroupStack}%
    \fi
  },
  % {<hue>}{<copies>}{<instances>}. The copies are drawn on a shifted canvas, which the bounding
  % box ignores, so the farthest one's corner is added to it.
  group stack/.style n args={3}{group copies={#1}{#2}, group count={#1}{#3},
    append after command={([shift={(#2*\diagramStackOffset, #2*\diagramStackOffset)}]
      \tikzlastnode.north east)}},
  % {<hue>}{<copies>}: from the farthest to the nearest, so the nearest lies on top.
  group copies/.code 2 args={%
    \ifnum#2>0
      \pgfmathtruncatemacro\diagramGroupNearerCopies{#2 - 1}%
      \edef\diagramGroupNearerCopies{{#1}{\diagramGroupNearerCopies}}%
      \pgfkeysalso{general shadow={subgroup=#1, shadow scale=1,
          shadow xshift=#2*\diagramStackOffset, shadow yshift=#2*\diagramStackOffset},
        group copies/.expand once=\diagramGroupNearerCopies}%
    \fi
  },
  % {<hue>}{<instances>}
  group count/.style 2 args={label={[mode chip=#1, anchor=north east, xshift=-\groupMargin,
    yshift=-\groupMargin]north east:$\times$\,#2}},
  % Titles with the same height and depth, so they line up whatever their letters.
  group title/.style={font=\diagramTitleFont\bfseries, text=#1, inner sep=0pt, text height=0.75em,
    text depth=0.2em},
  subgroup title/.style={font=\diagramNameFont\bfseries, text=#1, inner sep=0pt},
  info/.style={text=muted, align=left, inner sep=0pt},
  % Boxes: a state's lines are broken by hand, `box width` keeps the boxes of a column equal.
  box width/.style={minimum width=#1},
  state/.style={draw=#1, fill=white, line width=0.8pt, align=center,
    text=muted, inner xsep=1mm, inner ysep=2pt, rounded corners=\diagramCornerRadius},
  % A state at the figure's \boxWidth and \boxHeight.
  block/.style={state=#1, box width=\boxWidth, minimum height=\boxHeight},
  decision/.style={diamond, aspect=2, draw=#1, fill=white, line width=0.8pt, align=center,
    inner sep=0.5pt},
  % A state's mode, as a tab on its top right corner.
  mode chip/.style={fill=#1Tint, text=#1, inner xsep=2pt, inner ysep=1pt,
    anchor=south east},
  start/.style={circle, fill=ink, inner sep=0pt, minimum size=10pt},
  junction/.style={circle, fill=#1, inner sep=0pt, minimum size=3pt},
  junction/.default=wire,
  % Edges: a transition, a way back (dashed) and unplugging (red, dashed).
  bus/.style={draw=wire, line width=0.7pt},
  flow/.style={bus, ->},
  back/.style={flow, dash pattern=on 3pt off 2pt},
  cancel/.style={back, draw=unplug},
  edge label/.style={text=muted, align=center, inner sep=1.5pt},
  % Orthogonal routing: horizontal, vertical halfway between both ends, horizontal; or the vertical
  % at the x of a given coordinate (a bus).
  jog/.style={to path={-| ($(\tikztostart)!0.5!(\tikztotarget)$) |- (\tikztotarget) \tikztonodes}},
  jog via/.style={to path={-- (\tikztostart -| #1) |- (\tikztotarget) \tikztonodes}},
  % Circuits. A resistor lies on its wire, which is drawn first; a jack contact; a rail as a tag,
  % red for the pull-up rail, blue for ground.
  resistor/.style={draw=wire, fill=white, line width=0.6pt, minimum width=4mm,
    minimum height=1.5mm, inner sep=0pt},
  resistor vertical/.style={resistor, minimum width=1.5mm, minimum height=4mm},
  terminal/.style={circle, draw=wire, fill=white, line width=0.6pt, inner sep=0pt,
    minimum size=3pt},
  rail tag/.style={fill=#1!10, draw=#1, text=#1, line width=0.5pt,
    inner xsep=1.5pt, inner ysep=1pt, font=\diagramTextFont\bfseries,
    rounded corners=\diagramTagCornerRadius},
  % A plain area behind a part of a figure.
  backdrop/.style={fill=otherTint, rounded corners=\diagramCornerRadius},
  % The path a current or a reading takes, under the wires (on the background layer).
  highlight/.style={draw=#1!25, line width=3pt, line cap=round, line join=round},
  % A connection that only exists once a plug is inserted.
  link/.style={draw=other!70, line width=0.7pt, dash pattern=on 1pt off 1.5pt},
  % The mechanical link that sets a row of switches.
  actuator/.style={draw=muted, line width=0.4pt, dash pattern=on 2pt off 1pt},
  % A switch from its pivot at (0, -reach) to its contact at (0, reach); `yscale=-1` turns it
  % over. An open lever ends `switch reach` to the right, so the lever midpoints of a row of
  % switches, open or closed, lie on the line its actuator takes.
  switch reach/.store in=\diagramSwitchReach,
  switch reach=1.5mm,
  pics/switch/.style={code={
    \draw[bus] (0, -\diagramSwitchReach) -- (#1, \diagramSwitchReach);
  }},
  pics/switch open/.style={code={\pic {switch=\diagramSwitchReach};}},
  pics/switch closed/.style={code={\pic {switch=0pt};}},
}
